\chapter{Карта исходных материалов репозитория}
\label{app:source-map}

В этом приложении зафиксировано, какие материалы репозитория использованы в первом сборочном проходе ПЗ.

\begin{longtable}{|p{0.34\textwidth}|p{0.56\textwidth}|}
\caption{Связь материалов репозитория с разделами ПЗ}\label{tbl:source-map}\\
\hline
\textbf{Материал} & \textbf{Использование в ПЗ} \\
\hline
\endfirsthead
\hline
\textbf{Материал} & \textbf{Использование в ПЗ} \\
\hline
\endhead
\texttt{documentation/master-thesis-obsidian/thesis-draft/2/1} & Модель сопоставления резюме и вакансий, описание NLP-подхода. \\
\hline
\texttt{documentation/master-thesis-obsidian/thesis-draft/2/2} & Методика формирования тестового набора данных и план эксперимента. \\
\hline
\texttt{documentation/master-thesis-obsidian/thesis-draft/2/3} & Диаграммы AS IS и TO BE процесса рекрутмента. \\
\hline
\texttt{documentation/master-thesis-obsidian/thesis-draft/2/4} & Функциональные требования, Use Case и сценарий публикации вакансии. \\
\hline
\texttt{documentation/master-thesis-obsidian/thesis-draft/3/2} & Выбор и обоснование архитектуры сервиса. \\
\hline
\texttt{documentation/master-thesis-obsidian/thesis-draft/4/1} & Обоснование инкрементного жизненного цикла разработки. \\
\hline
\texttt{documentation/master-thesis-obsidian/thesis-draft/4/2} & Выбор IDE, языка программирования, фреймворков и хранилищ. \\
\hline
\texttt{version\_4/server.py}, \texttt{version\_4/client.py} & Текущая версия клиент-серверного прототипа, сохраненная в листингах. \\
\hline
\texttt{weaviate-quickstart/docker-compose.yml} & Конфигурационная заготовка для векторного хранилища. \\
\hline
\end{longtable}

\section{Листинги прототипа}

Полные листинги текущего прототипа оставлены отдельными файлами в каталоге \texttt{thesis/latex/listings}. При необходимости их можно включить в финальное приложение после согласования объема ПЗ.
